% ======================================================================
% CHAPTER 1: INTRODUCTION
% ======================================================================
\chapter{Introduction}

\section{Background of the Study}
The proliferation of diagnostic testing in modern healthcare has led to an unprecedented volume of medical reports being issued to patients. Blood panels, urinalysis, lipid profiles, thyroid function tests, radiological interpretations, and pathology reports are among the most commonly issued documents. Globally, billions of diagnostic tests are conducted annually, generating an enormous corpus of clinical data directed to patients and their healthcare providers. Despite their clinical importance, these reports are predominantly written in medical and scientific language that is inaccessible to the average patient. Reference ranges, clinical flags, and abbreviations are rarely self-explanatory, and the gap between report issuance and physician consultation can span days to weeks.

This information asymmetry has measurable consequences. Studies have indicated that patients who receive abnormal results without guidance frequently experience heightened anxiety or, conversely, fail to appreciate the significance of flagged values, leading to delayed follow-up care. Both outcomes are suboptimal from a public health perspective. A 2019 study published in the Journal of General Internal Medicine found that up to 40\% of patients receiving abnormal lab results experienced significant psychological distress while awaiting physician follow-up. While patient portals and digital health records have improved report delivery, they have done little to improve report comprehension, leaving a critical gap in the patient journey.

The emergence of large language models (LLMs) and natural language processing technologies has opened new possibilities for bridging this gap. Modern LLMs, trained on vast corpora including biomedical literature, can understand complex clinical text and generate coherent, accessible explanations for lay audiences. When combined with information retrieval techniques such as Retrieval-Augmented Generation (RAG), these models can ground their outputs in specific source documents, reducing the risk of hallucination and improving factual reliability. ReportRx leverages these capabilities to provide patients with immediate, context-aware interpretations of their medical reports, empowering them with knowledge while they await professional medical consultation.

Furthermore, the rapid digitization of healthcare records in India under the Ayushman Bharat Digital Mission has accelerated the availability of electronic medical reports. Combined with widespread smartphone penetration and affordable internet access, the technological infrastructure now exists to deploy AI-powered health literacy tools at scale. ReportRx is positioned to serve this emerging digital health ecosystem by providing an accessible, web-based platform that requires no specialised hardware or software installation. The application runs entirely in the browser, making it accessible across desktop and mobile devices, and requires no user registration or personal data storage, thus aligning with both user convenience and data privacy best practices.

\section{Problem Statement}
A large proportion of patients who receive medical reports lack the background to independently interpret their results. The delay between report receipt and professional consultation creates an unserved window during which patients may experience unnecessary anxiety or make uninformed decisions about their health. Existing general-purpose conversational AI tools, while capable of discussing medical topics, are not designed to systematically parse uploaded documents, identify out-of-range values, and produce structured, grounded analyses in a safe and disclaimered manner.

The key challenges that motivate this project are:

\begin{enumerate}
    \item \textbf{Lack of Patient-Facing Interpretation Tools:} Most existing healthcare informatics tools are built for clinicians and hospital administrators, not for the patient end-user. Patient portals typically display lab results as raw numerical values with reference ranges, placing the burden of interpretation entirely on the patient.
    \item \textbf{General-Purpose AI Limitations:} While general-purpose LLMs such as ChatGPT can answer medical questions, they are not designed for structured document analysis. They lack the ability to extract tabular data from PDFs, compare numerical values against reference ranges, and produce a consistent, structured report format with appropriate disclaimers.
    \item \textbf{Accuracy and Hallucination Risks:} Medical information is a high-stakes domain where inaccuracies can have serious consequences. Any system providing medical report interpretation must ground its outputs in the source document content and employ mechanisms to minimise hallucination, a challenge that standard chatbot architectures do not adequately address.
    \item \textbf{Privacy and Trust:} Medical documents contain highly sensitive personal information. Patients need assurance that their uploaded documents are handled confidentially and not stored or misused. Any practical solution must incorporate ephemeral document handling and clear communication about data practices.
    \item \textbf{Usability for Non-Technical Users:} The target audience includes patients of varying ages, educational backgrounds, and technological literacy levels. The system must offer an intuitive interface that requires no training to use, while still delivering comprehensive, accurate, and actionable information.
\end{enumerate}

ReportRx addresses these challenges by providing an AI-powered pipeline that: (i) accepts medical report uploads in PDF or image format; (ii) extracts and structures the textual content from these heterogeneous file formats; (iii) applies a dual-LLM architecture augmented by a RAG pipeline to produce a structured, contextually grounded analysis; and (iv) maps findings to evidence-based general health guidelines while clearly delineating the boundary between informational guidance and clinical advice.

\section{Objectives of the Project}
The primary objectives of the ReportRx project are outlined below:

\begin{enumerate}
    \item To design and build a web-based application that allows users to upload medical reports in PDF or image format and receive structured, lay-language analyses in real time, without requiring any account creation or personal data submission.
    \item To implement a dual-LLM architecture that combines a frontier general-purpose language model (Claude Sonnet / GPT-4) for fluent summarisation and explanation generation with a domain-specific biomedical model (BioMedLM) for accurate clinical terminology resolution and out-of-range value detection.
    \item To develop and deploy a Retrieval-Augmented Generation pipeline using LlamaIndex, ensuring that all model outputs are grounded in the actual content of the uploaded report and supplemented by a curated medical knowledge base, thereby minimising hallucination risk.
    \item To build a rule-based general health guideline engine that systematically maps identified out-of-range clinical parameters to standardised, evidence-based public health recommendations, and to ensure this mapping is deduplicated and contextually appropriate.
    \item To implement robust document processing capabilities supporting both PDF text extraction (via PyMuPDF) and OCR-based image extraction (via Tesseract), with appropriate error handling for unsupported or malformed file types.
    \item To design a responsive, accessible frontend interface that renders the analysis output in a clear, card-based layout with colour-coded flagging of abnormal values, collapsible sections, and a visually distinct, non-dismissible professional consultation disclaimer.
    \item To ensure responsible AI deployment through ephemeral document handling (no persistent storage of uploaded files), mandatory disclaimers on every output, and clear communication about the non-clinical, informational nature of the analysis.
    \item To evaluate the system's performance through systematic testing covering unit, integration, functional, and non-functional dimensions, with particular emphasis on response time, extraction accuracy, and output consistency.
\end{enumerate}

\section{Scope of the Project}

\subsection{In-Scope Elements}
ReportRx is scoped as an academic minor project with the following in-scope elements:

\begin{enumerate}
    \item \textbf{Supported Report Types:} Complete blood count (CBC), basic metabolic panel (BMP), comprehensive metabolic panel (CMP), lipid panel, thyroid function test (TFT), and standard urinalysis. These represent the most commonly ordered outpatient tests and cover a significant majority of routine diagnostic reports encountered in clinical practice.
    \item \textbf{Input Formats:} PDF documents (both text-based and scanned), and image files in JPG and PNG formats. The system differentiates between text-based PDFs (processed via PyMuPDF) and image-based PDFs (processed via Tesseract OCR).
    \item \textbf{Analysis Output:} The system produces a structured output comprising four sections --- (i) a lay-language summary of the entire report; (ii) key observations highlighting notable findings; (iii) a colour-coded table of out-of-range values with clinical context; and (iv) evidence-based general health guidance steps.
    \item \textbf{Language:} English-language reports only.
    \item \textbf{Deployment:} Web-based application accessible via modern web browsers on desktop and mobile devices.
\end{enumerate}

\subsection{Out-of-Scope Elements}
The following elements are explicitly out of scope for the current phase:

\begin{enumerate}
    \item \textbf{Clinical Diagnosis:} The system does not provide medical diagnoses, treatment recommendations, or medication advice under any circumstances.
    \item \textbf{Hospital Integration:} Real-time integration with hospital information systems (HIS) or electronic health record (EHR) platforms is not included.
    \item \textbf{Multi-Language Support:} Reports in languages other than English are not supported in this phase.
    \item \textbf{Data Persistence:} No patient data, uploaded reports, or analysis histories are stored. The system operates on an ephemeral, per-session basis.
    \item \textbf{Longitudinal Tracking:} Comparison of results across multiple reports over time is not supported.
    \item \textbf{Specialist Reports:} Radiology reports, pathology narratives, ECG interpretations, and other narrative-heavy or image-based specialist reports are not within scope.
\end{enumerate}

\subsection{Technology Scope}
The technology stack is confined to: Next.js (React) for the frontend; FastAPI (Python) for the backend; LlamaIndex for RAG orchestration; Claude Sonnet / GPT-4 and BioMedLM for LLM inference; PyMuPDF and Tesseract for document parsing; and Tailwind CSS for styling. No proprietary frameworks, paid licences, or enterprise tools are required beyond the LLM API access fees.
